% !TEX program = xelatex
% 공통수학2 · Ⅰ. 도형의 방정식 · 3. 도형의 이동 · 01 평행이동 · 1/1차시
% 학생용 활동지 (답안 없음)
\documentclass[11pt,a4paper]{article}
\usepackage{kotex}
\usepackage{iftex}
\ifPDFTeX\else
  \setmainhangulfont{Noto Serif CJK KR}
  \setsanshangulfont{Noto Sans CJK KR}
\fi
\usepackage[margin=2.1cm]{geometry}
\usepackage{parskip}
\usepackage{enumitem}
\usepackage{array}
\usepackage{tabularx}
\usepackage{xcolor}
\usepackage{amssymb}
\usepackage{tikz}
\usepackage[most]{tcolorbox}
\usepackage{fancyhdr}
\usepackage{titlesec}

\definecolor{goldc}{HTML}{B07C22}
\definecolor{goldstrong}{HTML}{8A5F16}
\definecolor{indigoc}{HTML}{2B3B57}
\definecolor{indigosoft}{HTML}{6E82A6}
\definecolor{linec}{HTML}{D6DACB}
\definecolor{surface2}{HTML}{E4E8DC}
\definecolor{writeline}{HTML}{C7CCBA}
\definecolor{inksoft}{HTML}{52604F}

\titleformat{\section}{\normalfont\sffamily\bfseries\large\color{indigoc}}{\thesection}{0.6em}{}
\titlespacing{\section}{0pt}{14pt}{6pt}
\newtcolorbox{goalbox}{enhanced, breakable, colback=white, colframe=goldc, boxrule=0.5pt, leftrule=3pt, arc=2pt, left=10pt, right=10pt, top=8pt, bottom=8pt}
\newtcolorbox{sheetbox}{enhanced, breakable, colback=white, colframe=linec, boxrule=0.6pt, arc=3pt, left=11pt, right=11pt, top=9pt, bottom=9pt}
\newcommand{\writespace}[1]{%
  \par\nobreak\vspace{3pt}
  \foreach \n in {1,...,#1} {%
    \noindent\textcolor{writeline}{\rule{\linewidth}{0.4pt}}\par\vspace{0.62cm}%
  }
}
\newcommand{\fillblank}[1]{\makebox[#1]{\hrulefill}}
\pagestyle{fancy}\fancyhf{}\renewcommand{\headrulewidth}{0pt}
\fancyfoot[L]{\footnotesize\color{inksoft} 공통수학2 · 2026학년도 2학기 · 지평선고등학교}
\fancyfoot[R]{\footnotesize\color{inksoft} 다음 시간 --- 02 대칭이동(1)}

\begin{document}

\noindent
\begin{minipage}[b]{0.62\linewidth}
  {\sffamily\footnotesize\color{indigosoft} 공통수학2 · Ⅰ. 도형의 방정식 · 3. 도형의 이동}\\[2pt]
  {\sffamily\bfseries\Large 01 평행이동 \textperiodcentered\ 14차시}
\end{minipage}\hfill
\begin{minipage}[b]{0.36\linewidth}
  \raggedleft\small\color{inksoft}
  반\ \fillblank{1.6cm} \quad 번호\ \fillblank{1.6cm}\\[4pt]
  이름\ \fillblank{3.6cm}
\end{minipage}

\vspace{4pt}
\noindent{\color{goldc}\rule{\linewidth}{2.2pt}}
\vspace{10pt}

\begin{goalbox}
\textbf{오늘의 목표}\\
점과 도형을 평행이동한 좌표·방정식을 구할 수 있다.
\vspace{4pt}
\small\color{inksoft}
$\square$ 자신 있음 \quad $\square$ 조금 헷갈림 \quad $\square$ 처음 봄 \hfill (수업 전 나의 상태에 표시)
\end{goalbox}

\section{생각 열기 --- 장기판 위의 이동}

\begin{sheetbox}
장기에서 차(車)는 한 수를 둘 때 선을 따라 한 방향으로 원하는 만큼 움직일 수 있다. 차를 A에서 B로 두 수 만에 옮기려면 어떻게 움직여야 할지 말해 보자. (가로/세로 각각 몇 칸씩?)
\writespace{3}
\end{sheetbox}

\section{개념 정리 --- 점과 도형의 평행이동}

\begin{sheetbox}
\textbf{점의 평행이동} --- 점 $(x,y)$를 $x$축의 방향으로 $a$만큼, $y$축의 방향으로 $b$만큼 평행이동한 점의 좌표는
\[
\left(\ \fillblank{2.4cm}\ ,\ \ \fillblank{2.4cm}\ \right)
\]

\vspace{8pt}
\textbf{도형의 평행이동} --- 방정식 $f(x,y)=0$이 나타내는 도형을 $x$축의 방향으로 $a$만큼, $y$축의 방향으로 $b$만큼 평행이동한 도형의 방정식은
\[
f(\ \fillblank{1.6cm}\ ,\ \ \fillblank{1.6cm}\ ) = 0
\]

\vspace{4pt}
{\footnotesize\color{inksoft} 주의! 점은 $+a, +b$를 더하는데, 도형의 방정식은 $x$ 대신 $x-a$, $y$ 대신 $y-b$를 대입한다. 왜 부호가 반대일까? 아래에 자신의 말로 적어 보자.}
\writespace{2}
\end{sheetbox}

\section{연습 문제}

\begin{sheetbox}
\textbf{1.} 다음 점을 $x$축의 방향으로 $-3$만큼, $y$축의 방향으로 $2$만큼 평행이동한 점의 좌표를 구하시오.\\
\hspace*{1.6em}⑴ $(-4,3)$ \qquad ⑵ $(5,-1)$
\writespace{2}

\vspace{6pt}
\textbf{2.} 원 $x^2+y^2=1$을 $x$축의 방향으로 $3$만큼, $y$축의 방향으로 $-1$만큼 평행이동한 원의 방정식을 구하시오.
\writespace{2}

\vspace{6pt}
\textbf{3.} 다음 방정식이 나타내는 도형을 $x$축의 방향으로 $-2$만큼, $y$축의 방향으로 $4$만큼 평행이동한 도형의 방정식을 구하시오.\\
\hspace*{1.6em}⑴ $x-4y-1=0$ \qquad ⑵ $(x+1)^2+(y-5)^2=9$
\writespace{4}

\vspace{6pt}
\textbf{4.} 이동식 수영장을 원 $(x-6)^2+(y-4)^2=2$로 나타낸다. 이 수영장을 $x$축의 방향으로 $a$만큼, $y$축의 방향으로 $b$만큼 평행이동하였더니 원 $x^2+y^2-2x-4y+3=0$이 되었다. 상수 $a$, $b$의 값을 구하시오.
\writespace{4}
\end{sheetbox}

\section{사고의 가시화 --- See · Think · Wonder}

\noindent{\footnotesize\color{inksoft} 오늘 배운 `평행이동'을 떠올리며 아래 세 칸을 채워 보자.}\\[6pt]
\begin{tabularx}{\linewidth}{@{}XXX@{}}
\textbf{\footnotesize\color{indigoc} See --- 무엇을 보았나?} &
\textbf{\footnotesize\color{indigoc} Think --- 무엇을 생각했나?} &
\textbf{\footnotesize\color{indigoc} Wonder --- 무엇이 궁금한가?} \\[4pt]
\end{tabularx}
\noindent
\begin{minipage}[t]{0.32\linewidth}\writespace{3}\end{minipage}\hfill
\begin{minipage}[t]{0.32\linewidth}\writespace{3}\end{minipage}\hfill
\begin{minipage}[t]{0.32\linewidth}\writespace{3}\end{minipage}

\section{마무리 성찰}

\begin{sheetbox}
\begin{minipage}[t]{0.48\linewidth}
{\footnotesize\color{inksoft} 오늘 배운 것을 한 문장으로}
\writespace{2}
\end{minipage}\hfill
\begin{minipage}[t]{0.48\linewidth}
{\footnotesize\color{inksoft} 감사 일기 / 유념 한 줄}
\writespace{2}
\end{minipage}
\end{sheetbox}

\end{document}
